Topologically close-packed intermetallics in the Fe--Mo system pose a severe data-efficiency problem because the technologically relevant complex phases contain many crystallographically distinct sites and large unit cells, making direct density-functional sampling prohibitively expensive. We therefore construct a machine-learning workflow that injects domain knowledge at three successive levels. Chemistry enters through Vegard-scaled prototype generation, crystallography enters through coordination-number-resolved averaging of atom-centered descriptors, and local bonding enters through physically motivated bond-order-potential moments together with ACE and SOAP reference representations. The curated repository records 292 usable Fe--Mo structures after curation, while the downstream error analysis scores 261 non-bcc/fcc/hcp structures with a phase-stratified 80/20 split and validates transfer on 70 additional complex-phase calculations comprising 19 $R$, 12 $P$, 11 $M$, and 28 $\delta$ structures. On the internal test split, the best archived kernel-ridge models reach root-mean-square errors of \SI{26.3}{\milli\electronvolt\per\atom} for ACE, \SI{28.2}{\milli\electronvolt\per\atom} for bond-specific BOP moments, and \SI{49.1}{\milli\electronvolt\per\atom} for SOAP, while coordination-number-resolved averaging lowers the BOP error by about 24\% and the SOAP error by about 40\% relative to the corresponding non-CNAV representations. The notebook validation workflow then propagates the trained models to the complex phases, where the stored parity plots show phase-resolved errors in the few-times-\SI{10}{\milli\electronvolt\per\atom} range and support subsequent Bragg--Williams calculations at \SI{1700}{\kelvin}. The thermodynamic analysis yields an $R$-phase mixing free-energy minimum of \SI{-53.4}{\milli\electronvolt\per\atom} near $x_{\mathrm{Fe}}=0.263$ and reproduces the measured evolution of site occupancies over the experimental composition interval $x_{\mathrm{Mo}}=0.27958$--0.37978. These results show that chemistry-aware structure generation, crystallography-aware aggregation, and bonding-aware descriptors together enable data-efficient prediction of complex transition-metal intermetallic energetics from a modest first-principles training set.
